\section{Một dụng cụ có kiểm định}
\label{sec:ch15-dung-cu-duoc-kiem}

\index{kiểm định chỉ số!ở họ mô hình khái niệm|(}%
Tiêu chí mà mục trước khép lại cho phép nói mô hình nào rò rỉ hơn mô hình nào,
và nó nói được như vậy vì người ta đã chấp nhận rằng hai điểm số ấy đo rò rỉ.
Mục này hỏi câu mà cả Phần~IV đã hỏi bảy lần và chưa lần nào nhận được câu trả
lời: lấy gì làm bằng chứng cho chỗ chấp nhận ấy. Bài báo đưa ba lập luận, và ba
lập luận ấy không cùng hạng.

\index{kiểm định chỉ số!lập luận thứ nhất, độ nhạy}%
Lập luận thứ nhất là \tn{độ nhạy}{sensitivity}, theo nghĩa thống kê của chữ ấy chứ không theo
nghĩa tiên đề mà chương~\ref{ch:attribution-theo-gradient} đã cấp cho nó: một
dụng cụ nhạy là dụng cụ báo khác nhau khi hai vật khác nhau.
Mục~\ref{sec:ch15-do-ro-ri} đã dùng lập luận này: cặp mô hình có sáu điểm đánh
giá lệch nhau không quá $0.005$ mà
\tn{điểm số can thiệp}{intervention score} lệch
$0.301$ được hai điểm số mới tách ra đúng chiều, trong khi OIS thì không, vì
khoảng tin cậy của OIS trên cặp ấy rộng tới mức chồng lên nhau. Lập luận này cho
thấy hai điểm số mới thấy được thứ mà các điểm đánh giá thông thường không thấy.

\index{kiểm định chỉ số!lập luận thứ hai và chỗ nó vòng lại}%
Lập luận thứ hai là hai điểm số \tn{rò rỉ}{leakage} tắt trên mô hình cứng. Bài báo huấn luyện mô
hình cứng trên tám bộ dữ liệu, ở cả chế độ khái niệm ít tương quan lẫn chế độ
khái niệm tương quan mạnh, và cả CTL lẫn ICL đều không lệch khỏi 0 một cách có ý
nghĩa ở bộ nào; OIS thì lệch ở cả tám. Lập luận này yếu hơn hai lập luận kia và
đáng nói rõ vì sao, vì chính Phần~IV đã dạy cách đọc kiểu này. Mô hình cứng được
tuyên là không rò rỉ \emph{theo định nghĩa}: huấn luyện độc lập tách hẳn việc học
khái niệm khỏi việc học nhãn, và phép nhị phân hóa trước khi vào đầu phân loại
ép giá trị học được về đúng giá trị thật, nên phần chênh của
định nghĩa~\ref{def:ch15-hai-diem-so} bằng 0 do cách dựng chứ không do quan sát.
Một điểm số báo khác 0 ở đó là một điểm số không khớp định nghĩa, và đó đúng là
điều đáng chê ở OIS, nhưng nó là phép kiểm \tn{tính nhất quán}{consistency} của
khung với chính nó
chứ không phải phép đối chiếu với thứ gì bên ngoài khung.

\index{rò rỉ!một nguồn chung bị đọc sai}%
Chữ \enquote{theo định nghĩa} ở trên chịu được sức nặng bao nhiêu thì có một
cách kiểm, và cách ấy dẫn ra ngoài danh mục. Cả bài của mục này lẫn khảo sát của
mục~\ref{sec:ch15-phan-con-lai} đều dẫn một bài báo thứ tư, bài của Mahinpei và
cộng sự về những hứa hẹn và cạm bẫy của mô hình học khái niệm
\tn{hộp đen}{black box}~\autocite{mclm}, và khảo sát nói về nó một điều mà đọc
bài gốc thì thấy ngược.

Khảo sát viết rằng bài ấy tìm ra biểu diễn cứng rò rỉ nhiều hơn hẳn biểu diễn
mềm. Phụ lục C của bài gốc viết ngược lại: khi thêm dần các khái niệm ngẫu nhiên,
tính dự báo của biểu diễn cứng có tăng, còn tính dự báo của biểu diễn mềm tăng
nhiều hơn, và \enquote{các biểu diễn mềm này luôn cho kết quả tốt hơn biểu diễn
cứng khi dùng trong nhiệm vụ phía sau}~\autocite{mclm}. Nhan đề mục 3 của bài
cũng nói vậy, rằng chỗ rò rỉ thông tin ngoại lai là \emph{trong biểu diễn mềm}.
Vậy câu ấy của khảo sát sai chiều, và mục~\ref{sec:ch15-phan-con-lai} sẽ trở lại
xem điều đó nói gì về một khảo sát.

\index{rò rỉ!hai đại lượng khác nhau cùng một tên}%
Nhưng bài gốc cũng không nói mô hình cứng an toàn, và đây mới là nửa có ích. Mục
4 của nó viết rằng ai đó có thể muốn dùng luôn biểu diễn cứng cho xong, rồi bác
chính ý ấy: ngay cả một số vừa phải những khái niệm cứng ngẫu nhiên và vô nghĩa
cũng nắm được nhiều thông tin về phân bố dữ liệu hơn người ta tưởng, nên
\enquote{rò rỉ thông tin có thể luôn là một vấn đề với các mô hình học khái niệm
hộp đen, bất kể dạng biểu diễn khái niệm}~\autocite{mclm}. Câu ấy và chữ
\enquote{theo định nghĩa} của bài mục này không mâu thuẫn nhau, vì chúng nói về
hai đại lượng khác nhau mang cùng một tên. Phần chênh của
định nghĩa~\ref{def:ch15-hai-diem-so} đo so với một tập khái niệm đã gán nhãn cố
định, nên với tập ấy giữ nguyên thì mô hình cứng cho phần chênh bằng 0. Bài gốc
thì thêm khái niệm vào rồi hỏi cả tập mang bao nhiêu thông tin về phân bố dữ
liệu, một câu hỏi mà phần chênh không đặt ra và không trả lời được. Sách giữ
phát biểu của bài, ghi rằng nó đúng trong phạm vi định nghĩa của chính nó, và
giữ nguyên đánh giá rằng lập luận thứ hai là phép kiểm nhất quán chứ không phải
phép kiểm định.

\index{kiểm định chỉ số!lập luận thứ ba, đối chiếu với điểm số can thiệp|(}%
Lập luận thứ ba mới là phép kiểm định, và nó là lập luận mà cả Phần~IV đi tìm.
Bài báo huấn luyện nhiều mô hình mềm và mô hình logit trên mỗi bộ dữ liệu, ở cả
hai chế độ tương quan và ở nhiều mức $\beta$, rồi đối chiếu ba điểm số với
điểm số can thiệp. Phép đối chiếu ấy phải đi vòng, vì mỗi điểm số của mỗi mô
hình không phải một con số mà là một con số kèm sai số: nó được đánh giá lại 5
lần trên mỗi mô hình, và năm lần ấy không cho cùng một giá trị. Bài xử lý chỗ đó
bằng cách coi mỗi điểm số là một phân bố chuẩn quanh giá trị trung bình của năm
lần, rút 10 nghìn mẫu từ phân bố ấy, và tính hệ số tương quan Pearson $r$ giữa
điểm số và điểm số can thiệp riêng cho từng mẫu. Kết quả là 10 nghìn giá trị $r$
thay vì một, và bước cuối gộp chúng lại thành một ước lượng kèm giá trị $p$ bằng
một quy tắc gộp có sẵn của thống kê, quy tắc Rubin, vốn được dựng để cộng đúng
hai nguồn sai số: sai số trong từng mẫu và sai số giữa các mẫu.
Bảng~\ref{tab:ch15-tuong-quan-diem-so} là kết quả.

\begin{table}[htbp]
  \centering
  \caption{Hệ số tương quan Pearson $r$ giữa ba điểm số rò rỉ và điểm số can thiệp,
  kèm giá trị $p$, trên tám bộ dữ liệu~\autocite{p30leakage}. Ba bộ đầu có khái
  niệm ít tương quan, ba bộ giữa có khái niệm tương quan mạnh, hai bộ cuối là dữ
  liệu thật. CTL vượt OIS ở cả tám dòng; ICL thì không, và ba dòng nó thua nằm ở
  bộ thứ hai, thứ năm và thứ sáu.}
  \label{tab:ch15-tuong-quan-diem-so}
  \begin{tabular}{@{}lcccccc@{}}
    \toprule
    & \multicolumn{2}{c}{CTL} & \multicolumn{2}{c}{ICL} & \multicolumn{2}{c}{OIS} \\
    \cmidrule(lr){2-3}\cmidrule(lr){4-5}\cmidrule(lr){6-7}
    Bộ dữ liệu & $r$ & $p$ & $r$ & $p$ & $r$ & $p$ \\
    \midrule
    TabularToy(0.25) & 0.75 & $3.4\times10^{-3}$ & 0.69 & $1.4\times10^{-2}$ & 0.49 & $1.2\times10^{-1}$ \\
    dSprites(0)      & 0.83 & $1.1\times10^{-3}$ & 0.68 & $1.5\times10^{-2}$ & 0.70 & $1.1\times10^{-2}$ \\
    3dshapes(0)      & 0.84 & $4.2\times10^{-5}$ & 0.86 & $2.4\times10^{-5}$ & 0.82 & $2.0\times10^{-4}$ \\
    \addlinespace
    TabularToy(0.75) & 0.54 & $4.9\times10^{-2}$ & 0.51 & $1.1\times10^{-1}$ & 0.27 & $4.1\times10^{-1}$ \\
    dSprites(4)      & 0.67 & $1.9\times10^{-2}$ & 0.48 & $1.2\times10^{-1}$ & 0.55 & $7.4\times10^{-2}$ \\
    3dshapes(5)      & 0.79 & $3.9\times10^{-4}$ & 0.69 & $3.5\times10^{-2}$ & 0.70 & $6.5\times10^{-3}$ \\
    \addlinespace
    CUB              & 0.72 & $9.0\times10^{-3}$ & 0.67 & $2.1\times10^{-2}$ & 0.44 & $1.1\times10^{-1}$ \\
    HAM10K           & 0.49 & $1.8\times10^{-2}$ & 0.41 & $5.4\times10^{-2}$ & $-0.06$ & $7.7\times10^{-1}$ \\
    \bottomrule
  \end{tabular}
\end{table}

\index{kiểm định chỉ số!chỗ bài báo nói quá bảng của mình}%
Bảng ấy nói được một điều mạnh và một điều yếu hơn điều mà bài báo viết. Điều
mạnh là CTL: nó tương quan với điểm số can thiệp cao hơn OIS ở cả tám bộ, kể cả ở
HAM10K nơi OIS xuống $-0.06$, và giá trị $p$ của nó dưới $0.05$ ở cả tám, chỗ sát
nhất là TabularToy(0.75) với $4.9\times10^{-2}$. Điều yếu hơn là ICL. Bài báo
viết rằng hai điểm số \enquote{vượt OIS một cách hệ thống trên mọi bộ dữ liệu},
và bảng của chính nó bác câu ấy ở ba dòng: dSprites(0) cho ICL $0.69$ so với OIS
$0.70$, dSprites(4) cho $0.48$ so với $0.55$, và 3dshapes(5) cho $0.69$ so với
$0.70$. Đọc theo giá trị $p$ thì ICL cũng thua ở đúng ba dòng ấy. Sách giữ phát
biểu cho CTL, bỏ phần về ICL, và không dùng ICL làm bằng chứng ở bất cứ đâu
trong chương. Không kết luận nào của bài phụ thuộc vào chỗ ấy: hai câu ngay sau
đó trong chính bài đã phát biểu đúng phần yếu hơn, rằng CTL thường tương quan
mạnh hơn ICL và có ý nghĩa ở mọi bộ dữ liệu.

\index{kiểm định chỉ số!lập luận thứ ba, đối chiếu với điểm số can thiệp|)}%
\index{ground truth!đường thứ tư, gán nhãn sẵn trong dữ liệu}%
Điều làm phép kiểm định ấy chạy được không nằm trong khung thông tin nào cả, mà
nằm ở chỗ bộ dữ liệu có gán nhãn khái niệm. Sách đã đi qua ba đường dựng ground
truth và mục~\ref{sec:ch14-ground-truth-ke-tan-cong} đã đếm chúng: thiết kế
nhiệm vụ ở chương~\ref{ch:chi-so-khong-do-do-trung-thuc}, thiết kế dữ liệu, và
thiết kế mô hình ở chương~\ref{ch:giai-thich-duoi-tan-cong}. Đây là đường thứ
tư, và nó là đường rẻ nhất trong bốn, vì không ai phải dựng gì cả: khi bộ dữ
liệu đã có 112 thuộc tính gán tay cho mỗi bức ảnh chim, thì lượng thông tin mà
mỗi thuộc tính mang về loài chim là một đại lượng tính thẳng từ bảng nhãn. Cả hai
mốc trong định nghĩa~\ref{def:ch15-hai-diem-so} đến từ chỗ ấy, và điểm số can thiệp
cũng vậy.

\index{kiểm định chỉ số!ba chỗ phép kiểm định không với tới}%
Ba chỗ mà phép kiểm định ấy không với tới, và chương nêu cả ba vì bỏ chỗ nào thì
mục này thành lời khen. Thứ nhất, tiêu chí dùng để kiểm định không phải một tiêu
chí tốt, chỉ là một tiêu chí độc lập. Mục~\ref{sec:ch15-do-ro-ri} đã liệt các
chỗ hỏng của điểm số can thiệp theo lời chính bài báo, và bốn chỗ ấy không mất đi
khi nó được dùng làm mốc: một dụng cụ đối chiếu tốt với một tiêu chí đặc hiệu mà
không nhạy thì kế thừa đúng tính đặc hiệu ấy. Nói được là hai điểm số bắt được
những gì điểm số can thiệp bắt được, và không nói được gì về những gì
điểm số can thiệp bỏ sót.

Thứ hai, phép ước lượng thông tin tương hỗ có độ chệch phụ thuộc số chiều của
biến, nên hai đại lượng ước lượng trên hai biểu diễn khác số chiều thì không trừ
nhau được. Với mô hình mà giá trị khái niệm là một số như ba loại của
mục~\ref{sec:ch15-khai-niem-trong-mo-hinh} thì không sao, vì cả giá trị học được
lẫn nhãn thật đều là một số. Với kiến trúc mà mỗi khái niệm mang cả một vector,
lớp mô hình mà mục~\ref{sec:ch15-nguyen-nhan} đọc, thì phép trừ trong
định nghĩa~\ref{def:ch15-hai-diem-so} không thực hiện được, và bài báo phải bỏ
mốc đi, giữ lại vế đầu. Cái còn lại sau khi bỏ mốc không phải điểm số rò rỉ mà
là lượng thông tin thô, so được giữa các mô hình cùng dạng và không so được với
0.

Thứ ba là ranh giới mà phần mở đầu chương đã nêu và đây là chỗ nó cụ thể ra.
Thứ được kiểm định là một điểm số rò rỉ, tức một điểm số đo cột $\hat c_k$ mang
thêm bao nhiêu so với một tập khái niệm đã gán nhãn. Nó không phải một chỉ số độ
trung thực, và bảy chương của Phần~IV vẫn ở nguyên chỗ chúng đứng: không có
dòng nào trong bảng~\ref{tab:ch15-tuong-quan-diem-so} nói gì về deletion, về
insertion, hay về việc một điểm số attribution có nêu đúng đặc trưng mà mô hình
dựa vào hay không. Điều chương này thêm vào là một mẫu, rằng có tồn tại một hoàn
cảnh trong ngành nơi một dụng cụ đo được đối chiếu với một tiêu chí độc lập, và
hoàn cảnh ấy có một đặc điểm nhận dạng được: ground truth nằm sẵn trong bộ dữ
liệu chứ không phải do ai dựng.

\index{kiểm định chỉ số!ở họ mô hình khái niệm|)}%
Đặc điểm ấy là thứ chương~\ref{ch:ground-truth-dung-san} nhận, và nhận kèm câu
hỏi hiển nhiên đi theo: một bộ dữ liệu có gán nhãn khái niệm thì đắt tới mức
nào, và khi tập khái niệm ấy thiếu thì chuyện gì xảy ra. Mục sau trả lời vế thứ
hai.
